% Generated draft: requires mathematical and source review.
\documentclass[11pt]{article}
\usepackage[margin=25mm]{geometry}
\usepackage{fontspec}
\setmainfont{Latin Modern Roman}
\usepackage{amsmath,amssymb,mathtools,unicode-math}
\setmathfont{Latin Modern Math}
\usepackage{enumitem,needspace,xurl,hyperref}
\setlength{\parindent}{0pt}
\setlength{\parskip}{5pt}
\setlength{\emergencystretch}{3em}
\newcommand{\N}{\mathbb{N}}
\newcommand{\Z}{\mathbb{Z}}
\newcommand{\Q}{\mathbb{Q}}
\newcommand{\R}{\mathbb{R}}
\newcommand{\C}{\mathbb{C}}
\newcommand{\F}{\mathbb{F}}
\newcommand{\A}{\mathbb{A}}
\newcommand{\PP}{\mathbb{P}}
\newcommand{\E}{\mathbb{E}}
\newcommand{\eps}{\varepsilon}
\newcommand{\dd}{\,\mathrm{d}}
\begin{document}
\section*{UnsolvedMath ID: 3100040}
\textbf{Source problem code:} catalogue entry
\subsection{Definitions and mathematical statement}
For an integer $n\geq1$, let $P_n$ be the set of integer partitions $\lambda=(\lambda_1,\ldots,\lambda_r)$ with $\lambda_1\geq\cdots\geq\lambda_r>0$ and $\sum_i\lambda_i=n$. Say $\lambda\preceq\mu$ if the parts of $\lambda$ can be grouped into disjoint nonempty blocks whose sums are the parts of $\mu$; equal-valued parts may occur and are grouped by their indices. This is refinement, not dominance order or inclusion of Ferrers diagrams.

The poset is graded by $\operatorname{rk}(\lambda)=n-r$: merging two parts increases rank by one. A set of pairwise incomparable elements is an antichain. A finite ranked poset is Sperner when its largest antichain has the same cardinality as its largest rank level.

\textbf{Question.} Is $P_n$ Sperner for every integer $n\geq1$? Equivalently, if $p_r(n)$ is the number of integer partitions of $n$ into exactly $r$ parts, is
\[
\max\{|\mathcal A|:\mathcal A\subseteq P_n\text{ is an antichain}\}
=\max_{1\leq r\leq n}p_r(n)
\]
for every $n$? The source does not explicitly mention $n$; the displayed statement is the conventional fixed-sum family interpretation, with the universal quantifier made explicit. An unrestricted union of partitions of all integers would instead be an infinite disjoint union of these refinement posets.

The supplied research summary mentions possible small-case progress but explicitly expresses uncertainty and supplies no verified theorem or data. Its discussion of the set-partition lattice is about a different poset and does not settle this integer-partition question.
\subsection{Short English statement}
For every fixed integer sum, are the largest incomparable families of integer partitions no larger than the largest family with a fixed number of parts? The source concerns refinement of integer partitions, not labelled set partitions.
\subsection{Sources}
\begin{itemize}
\item Cooper, \emph{Combinatorial Problems I Like}, Matroids, Lattices, and Posets, source-order bullet 40, October 2020 snapshot, \url{https://people.math.sc.edu/cooper/combprob.html}; HTML extraction, accessed 2026-07-29. Later resolutions were not checked item by item according to the record.
\item The supplied summary speculates about small cases and related partition lattices but identifies no publication, theorem, or verified computation for this refinement poset.
\item Local registry \texttt{data/unsolvedmath/problems.json}, permanent ID $3100040$, source code catalogue entry, \emph{Is the poset of integer partitions ordered by refinement Sperner}. The record labels this partially solved and PARTIAL-PROGRESS while its summary says likely open, with NEEDS\_REVIEW source-status and rights qualifications.
\end{itemize}
\end{document}
